% Post-CP21 candidate; exact universal polynomial certificate retained.
\subsection{Degree-one full image on every principal surface of two-rank zero}
\status{Proved and independently checked.}

\begin{lemma}[The polar tensor for every genus-two curve of two-rank zero]
\label{lem:p-rank-zero-genus-two-polar-stable-candidate}
Let $C$ be a smooth genus-two curve over an algebraically closed field
of characteristic two, and suppose its Jacobian has two-rank zero.
In the complete three-theta embedding of its Jacobian, the alternating
polar three-form of its Coble cubic is stable for $\operatorname{SL}_9$.
\end{lemma}
\begin{proof}
We first make the range of normal forms explicit. A genus-two curve
is hyperelliptic, and its degree-two canonical map is separable:
a purely inseparable degree-two extension of $k(x)$ would be
$k(x^{1/2})$, and hence rational. Choose a geometric Weierstrass
point as infinity. Its equation has the form
\[
       y^2+h(x)y=f_5(x),\qquad
       h=h_0+h_1x+h_2x^2,
\]
where $f_5$ has degree five and $h\ne0$. This is also the curve
normal form in \citet[Proposition~2.1]{RainsSam2018Trivectors}.
A basis of regular differentials is $\omega=dx/h$ and $x\omega$.
The rational Cartier rule gives
\[
 C\omega=\sqrt{h_1}\,\omega,
 \qquad C(x\omega)=(\sqrt{h_0}+\sqrt{h_2}x)\omega.
\]
Indeed one writes $dx/h=h\,dx/h^2$ and similarly for $x\,dx/h$;
Cartier retains exactly the odd powers of $x$. This triangular
semilinear matrix has stable rank zero precisely when
$h_1=h_2=0$. Thus $h=h_0\ne0$. Scaling $y$, then $x$, normalizes
the equation to $y^2+y=x^5+$ lower terms. Changes $y\mapsto y+u(x)$
remove the even powers and constant term over the algebraically
closed field. In particular every curve under consideration occurs
in the slightly larger family
\begin{equation}\label{eq:p-rank-zero-normal-form-candidate}
                  y^2+y=x^5+a x^3+b x^2+c x.
\end{equation}
Every member of this family is smooth, including its unique point
at infinity: its affine $y$-derivative is one, and its odd pole of
order five gives the smooth genus-two Artin--Schreier model at infinity.

Use the normal form of the alternating tensor in the same cited
Proposition~2.1, with $c_{15}=1,c_{12}=a,c_{18}=b,c_{24}=c$ and
all other coefficients zero. With $[ijk]=e_i\wedge e_j\wedge e_k$,
it is
\[
 \gamma=[267]+[258]+[348]+[169]+[357]+[249]+[178]+[456]
                 +[235]+a[147]+b[145]+c[134].
\]
\citet[Proposition~2.13 and Proposition~3.8, including their proofs]
{RainsSam2018Trivectors} show that this tensor is stable, that its
rank-four surface is the Jacobian of \eqref{eq:p-rank-zero-normal-form-candidate}
in its complete three-theta embedding, and that its Pfaffian cubic
is the corresponding Coble cubic. Let $B$ be the alternating polar
three-form of that cubic. The following explicit expression and
minor identity are over the polynomial ring $\mathbf F_2[a,b,c]$:
\begin{align*}
 B={}&[129]+[138]+[156]+[179]+[236]+[259]+[268]
                 +[349]+[367]+[457]\\
    &+a\bigl([134]+[148]+[157]\bigr)
                 +b[146]+c\bigl([126]+[135]\bigr).
\end{align*}
To specify the minor without any numerical specialization, form
\[
 \mathcal D_B:\operatorname{End}(k^9)\longrightarrow\bigwedge^3 k^9,
                   \qquad X\longmapsto X\cdot B.
\]
Its 81 columns are $E_{ij}$ in lexicographic order and its 84 rows
are the lexicographically ordered triples. Delete the three rows
$[467],[478],[578]$. The resulting square matrix satisfies
\begin{equation}\label{eq:p-rank-zero-polar-unit-minor-candidate}
                         \det\mathcal D'_B=1
                   \quad\text{in }\mathbf F_2[a,b,c].
\end{equation}
This is a finite polynomial identity, with a retained exact
unit-pivot elimination certificate described below. It proves that
the infinitesimal stabilizer of $B$ in $\operatorname{GL}_9$ is zero
for every value of $a,b,c$, and hence also in $\operatorname{SL}_9$.

It remains to justify why this infinitesimal test proves stability
in this setting. Since $\gamma$ is stable, it is conjugate into the
standard Cartan subspace, by
\citet[Corollary~5.6]{RainsSam2018Trivectors}.
The Pfaffian cubic construction and ordinary cubic polarization are
covariant under change of basis, up to a nonzero determinant scalar.
In a Cartan basis the cubic is the sum of pure cubes and affine-line
products, so its alternating polar tensor is again in the same
four-dimensional Cartan subspace. Therefore $B$ is conjugate, up
to scalar, to a Cartan tensor. By the same source's
Propositions~2.9--2.10, a nonstable Cartan tensor lies on a reflection
hyperplane and has a positive-dimensional stabilizer. Such a
stabilizer has nonzero Lie algebra. This contradicts
\eqref{eq:p-rank-zero-polar-unit-minor-candidate}. Thus $B$ is stable.
\end{proof}

\begin{remark}[A universal polynomial certificate]
\label{rem:p-rank-zero-polar-unit-certificate-candidate}
The executed exact program is
\path{agents/monodromy/code/ppav_supersingular_polar_uniform_lie.py},
with full output
\path{agents/monodromy/outputs/ppav_supersingular_polar_uniform_lie.json}.
It builds the alternating contraction matrix directly from the twelve
terms of $\gamma$, computes all nine principal eight-Pfaffians by
recursive matching, and verifies their common cubic after division
by the omitted coordinate. It then obtains $B$ from the squarefree
cubic monomials and constructs the infinitesimal action matrix from
$E_{ij}\cdot[e_r\wedge e_s\wedge e_t]$ by the exterior-product rule.
Polynomial Gaussian elimination uses only pivots equal to the constant
one; all 81 pivot pairs and all coefficients are retained.
The three unused exterior rows are exactly $[467],[478],[578]$;
the added trace row is unused as well. Thus the elimination proves
the displayed square-minor identity over the entire polynomial ring,
not just at a list of field values. The author execution at
18:29:35.876074 UTC on 19 November 2026 took
0.019472674997814465 seconds. The parent, reduction/descent, and obstruction/crystal agents independently
audited every proof and reread the cited original sources. The parent
independently reconstructed the displayed tensor action and the fixed minor
in 0.06436865900468547 seconds; reduction/descent independently rebuilt
all nine Pfaffians and that minor in 0.10397645600460237 seconds;
obstruction/crystals independently reconstructed the displayed minor
in 0.06390757599729113 seconds. Their author-output-preserving executions
of the complete retained program took respectively
0.02088291999825742, 0.02283291499770712, and
0.02654484500089893 seconds. All checks passed.
The earlier finite normal-form and Lie probes remain retained as
discovery evidence; they are not the proof of the uniform identity.
\end{remark}

\begin{theorem}[Every fixed principal surface of two-rank zero]
\label{thm:p-rank-zero-ppav-degree-one-full-candidate}
Let $(A,\Theta)$ be a principally polarized abelian surface of
geometric two-rank zero over any perfect field $k$ of characteristic
two. For the smooth maximal-rank locus in $|3\Theta|$, the relative
unit local system has rank eight, and both its geometric and
arithmetic images are $\GL_8(\mathbf F_2)$.
In particular the representation is geometrically irreducible.
\end{theorem}
\begin{proof}
Work first over an algebraic closure. If the principal surface is
a product, both elliptic factors are supersingular. In characteristic
two they are geometrically isomorphic to $y^2+y=x^3$, and the result
is Theorem~\ref{thm:ss-product-degree-one-full-group}.
Otherwise $A$ is the Jacobian of its smooth genus-two theta divisor.
The same is true of $A^{(2)}$, whose two-rank remains zero.
The preceding lemma proves stability of the alternating polar of
the Coble cubic of $A^{(2)}$.

Lemma~\ref{lem:root-nonordinary-Coble-identification-candidate}
identifies the actual Cartier Pfaffian cubic with this Coble cubic
in the exact complete three-theta embedding. Consequently their
polar tensors are identified under the same projective linear
change of basis and a nonzero scalar. No identification of the
actual Cartier trivector with a particular marked-curve trivector
is needed for this step: equality of the cubics is sufficient.
If different principal line-bundle representatives are used,
translation on the abelian surface identifies the complete systems
and the associated embeddings. Finally stability is unchanged by
the coefficient-square-root automorphism of the algebraically
closed field. The extra hypothesis in
Theorem~\ref{thm:nonordinary-ppav-polar-stable-conditional-candidate}
is therefore satisfied, and that theorem gives the full geometric
image. The arithmetic image contains it and lies in the same
$\GL_8(\mathbf F_2)$, giving equality over the original perfect field.
\end{proof}

\begin{remark}[The remaining degree and two-rank-one gates]
\label{rem:p-rank-zero-ppav-degree-one-scope-candidate}
This theorem concerns degree one for the complete three-theta
embedding. It does not assert higher-degree full images on a fixed
nonproduct surface, or degree-one full images on every fixed
nonproduct surface of two-rank one. Those are separate structural
gates. The earlier conditional polar theorem is nonvacuous on the
actual Cartier locus by the preceding uniform result, while its
extra hypothesis still requires investigation in two-rank one.
\end{remark}
